\section{Analysis: Why Memory Helps or Hurts}
\label{sec:analysis}

\textbf{Failure taxonomy (baseline).} Over the \Neval{}-run M0 snapshot, a systematic per-run triage (machine-checkable evidence, explicit UNKNOWN share, never folding unclassified runs) yields: (i) no-output/aborted 21.0\% (146/696), (ii) semantic/partial failure 62.8\% (437/696; 132 runtime crashes + 305 assertion failures), (iii) contract/format 0.3\% (2/696), residual UNKNOWN 7.0\% (49/696). The dominance of semantic failures (62.8\%)---deliverables that are syntactically plausible but semantically wrong---is direct evidence that the binding constraint is \emph{knowledge of what to do}, not output formatting. This supports the memory hypothesis: an agent that consolidated ``which assertions fail and why'' across sessions could convert semantic failures into passes.

\textbf{Trace anatomy (baseline).} 47 instrumented pilot runs (2026-11-03--06) show: median 26.5 tool calls (IQR 13.8--38.0, max 97) across 20.5 iterations, median 1{,}047\,s (IQR 552--1{,}495). Of 1{,}371 tool calls, 41.2\% go through the MCP bridge to live EDA tools and 58.8\% are built-in file/search/shell. The effort is substantial but not unbounded---and the 41.2\% MCP fraction means nearly half of all actions require grounding in live tool outputs the model has never seen during training. (Caveat: these 47 traces are archived pilots, not the \Neval{}-run batch; they describe mechanics, not pass rates.)

\textbf{Affective-mechanism analysis (placeholder, post pre-registration).} The three mechanisms (\S\ref{sec:affective}) predict testable signatures: (1) surprise distribution---M5 should concentrate writes on high-surprise events; (2) $r_{\text{cp}}$ vs.\ Pass@1 scatter---M7 should occupy the ``smaller memory, equal-or-better Pass@1'' region; (3) state\_match hit distribution---M6 should up-weight pitfall memories in the fail-streak subset; (4) M4 leak-rate audit---transfer-hit memories should be abstract cognition, not problem-specific answer fragments (leak rate $=0$). These analyses are \TBD{} pending post-pre-registration evaluation.

\textbf{Why memory might hurt.} If memory is uncurated (M2, no reliability gating), stale or contradictory entries can mislead: a lesson from an old tool version may be wrong for the current one. This is why H1a tests M3 (gated) against M2 (ungated), and why the disconfirmation filter and held-out rollback are part of the reliability gating. If H1a falsifies \emph{but} M3 $\approx$ M2, the negative result itself is publishable: it bounds the value of reliability gating in long-chain flow.
